\documentclass[10pt,a4paper]{amsart}
\usepackage[margin=19mm]{geometry}
\usepackage{amsmath,amssymb,mathtools}
\usepackage[scr=rsfs,cal=euler]{mathalfa}
\usepackage{xfrac}
\usepackage[unicode,hidelinks,bookmarks=false]{hyperref}
\newcommand{\ie}{i.e.\ }
\newcommand{\cf}{cf.\ }
\newcommand{\resp}{respectively\ }
\newcommand{\Subsection}{\subsection}
\providecommand{\nobreakdash}{\nobreak\hbox{-}}
\setlength{\emergencystretch}{2em}
\multlinegap0pt
\usepackage{amssymb}
\usepackage{tikz}
\newtheorem{definition}{Definition}
\newtheorem{lemma}{Lemma}
\newcommand{\R}{\mathbb R}
\newcommand{\abs}[1]{\lvert #1\rvert}
\newcommand{\dist}{\operatorname{dist}}
\title{Rate of convergence of the $p$-torsion function to the distance function}
\author{Farid Bozorgnia, Masoud Bayrami-Aminlouee, Khudoyor Mamayusupov}
\date{Source: arXiv:2609.06534v1 (CC BY 4.0)}
\begin{document}
\maketitle

\noindent\emph{Source and licence.} Rate of convergence of the $p$-torsion function to the distance function. Version arXiv:2609.06534v1 (CC BY 4.0), distributed by arXiv under the Creative Commons Attribution 4.0 International licence, http://creativecommons.org/licenses/by/4.0/, as recorded in the arXiv licence metadata for this exact version and shown on the abstract page https://arxiv.org/abs/2609.06534v1. Full licence notice: \texttt{licenses/arxiv-2609.06534-CC-BY-4.0.txt}.\par\smallskip

\noindent\textit{Editorial scope.} Counted unit: the article's lemma lem:tv-exterior-ball (source0.tex lines 408--416). The article's complete original proof of it is delivered with it (source0.tex lines 417--533, one \texttt{proof} environment of 117 lines). No mathematics is written, repaired or re-derived: the statement and the proof are the article's own lines, in the article's own notation, and no retained line is substituted or re-flowed.  Retained uncounted as the source-local context the proof needs: lines 370--386; lines 400--402.  Nothing was omitted from any retained span, so the package carries no omission record.  No retained line contains a citation, so the wrapper carries no bibliography and no bibliography entry.  No retained line contains a figure environment, an image-inclusion command or drawing code, so no image is delivered and the package declares no asset dependency.  Editorial formatting only, in addition: an AMS \texttt{amsart} preamble that restates the article's own declarations and macros used by the retained lines, namely the environments \texttt{definition}, \texttt{lemma} on separate counters, together with the standard packages those lines need.  The retained environments print in this document as Definition 1, Lemma 1 in order of appearance, whereas the article prints the same items with the numbering of its own full document.  The complete native source of the article as distributed in the arXiv e-print for this exact version is bundled unchanged as \texttt{sources/209544/source0.tex}.
\begin{definition}[Uniform exterior ball condition]\label{def:ext-ball}
A bounded domain $\Omega\subset\R^n$ satisfies the \emph{uniform exterior ball
condition with radius $r_e>0$} if, for every $y\in\partial\Omega$, there
exists $c_y\in\R^n$ such that
\[
\abs{c_y-y}=r_e,
\qquad B(c_y,r_e)\cap\Omega=\emptyset,
\]
where $B(c_y,r_e)$ denotes the open ball; $c_y$ is not assumed to be unique
and no boundary regularity is assumed in this definition. If
$\partial\Omega$ is of class $C^1$ with outer unit normal $\nu$, then
necessarily $c_y=y+r_e\nu(y)$, and the condition takes the familiar form
\begin{equation}\label{eq:ext-ball}
B\bigl(y+r_e\nu(y),r_e\bigr)\cap\Omega=\emptyset
\qquad\text{for every }y\in\partial\Omega .
\end{equation}
\end{definition}

\begin{equation}\label{eq:tangency}
\abs{x-c_y}=r_e+d(x).
\end{equation}

\begin{lemma}[Finite total variation of the distance Laplacian]\label{lem:tv-exterior-ball}
Let $\Omega\subset\R^n$ be a bounded $C^1$ domain satisfying
Definition~\ref{def:ext-ball} with radius $r_e>0$. Then $\mu:=-\Delta d$ is a
signed Radon measure on $\Omega$ with
\[
\abs{\mu}(\Omega)\ \le\ 2c_\Omega+\frac{2n}{r_e}\abs{\Omega}\ <\ \infty,
\]
where $c_\Omega$ is the boundary-layer constant of \eqref{eq:layer} below.
\end{lemma}

\begin{proof}
\emph{Step 1: $d$ is a minimum of smooth functions.}
For $y\in\partial\Omega$ let $c_y$ be as in Definition~\ref{def:ext-ball} and
put $g_y(x):=\abs{x-c_y}-r_e$. Fix $x\in\Omega$. Since $\abs{x-c_y}>r_e$,
$g_y(x)$ is the distance from $x$ to the ball $B(c_y,r_e)$, and since that
ball is contained in $\R^n\setminus\Omega$,
\[
g_y(x)=\dist\bigl(x,B(c_y,r_e)\bigr)\ \ge\ \dist(x,\R^n\setminus\Omega)=d(x)
\qquad\text{for every }y\in\partial\Omega,
\]
with equality when $y$ is a nearest boundary point of $x$, by the tangency
identity \eqref{eq:tangency}. Hence
\begin{equation}\label{eq:d-as-min}
d(x)=\min_{y\in\partial\Omega}g_y(x)\qquad(x\in\Omega).
\end{equation}

\emph{Step 2: semiconcavity and the sign of $\mu+\frac n{r_e}\mathcal L^n$.}
Each $g_y$ is smooth on $\Omega$ with
\[
D^2g_y(x)=\frac{1}{\abs{x-c_y}}\bigl(\mathrm{Id}-\hat n\,\hat n^{\top}\bigr)
\ \le\ \frac1{r_e}\,\mathrm{Id},
\qquad \hat n:=\frac{x-c_y}{\abs{x-c_y}},
\]
because $\abs{x-c_y}>r_e$. Thus on every ball $B\Subset\Omega$ the functions
$g_y-\tfrac1{2r_e}\abs{x}^2$ are concave, and by \eqref{eq:d-as-min} so is
their infimum $f:=d-\tfrac1{2r_e}\abs{x}^2$. A function that is concave on
every ball $B\Subset\Omega$ satisfies $\partial_{\xi\xi}f\le0$ in
$\mathcal D'(\Omega)$ for every unit vector $\xi$: for non-negative
$\varphi\in C_c^\infty(B)$ and $h$ smaller than the distance from
$\operatorname{supp}\varphi$ to $\partial B$,
\[
\int_\Omega f\,\partial_{\xi\xi}\varphi\,dx
=\lim_{h\to0}\int_\Omega
\frac{f(x+h\xi)-2f(x)+f(x-h\xi)}{h^{2}}\,\varphi(x)\,dx\ \le\ 0,
\]
since the second difference of a concave function is non-positive, and a
partition of unity extends this to every non-negative
$\varphi\in C_c^\infty(\Omega)$. As $\partial_{\xi\xi}\bigl(\tfrac1{2r_e}\abs{x}^2\bigr)=r_e^{-1}$
for $\abs\xi=1$, this gives
\[
\langle\partial_{\xi\xi}d,\varphi\rangle\ \le\ \frac1{r_e}\int_\Omega\varphi\,dx
\qquad\text{for every unit }\xi\text{ and every non-negative }
\varphi\in C_c^\infty(\Omega),
\]
which is the meaning of $D^2d\le r_e^{-1}\mathrm{Id}$ in $\mathcal D'(\Omega)$.
Taking the trace, that is, summing over an orthonormal basis $e_1,\dots,e_n$,
yields $\langle\Delta d,\varphi\rangle\le\frac n{r_e}\int_\Omega\varphi\,dx$,
or equivalently
\[
\Bigl\langle\mu+\frac n{r_e}\,\mathcal L^n,\ \varphi\Bigr\rangle\ \ge\ 0
\qquad\text{for every non-negative }\varphi\in C_c^\infty(\Omega).
\]
Hence $\widetilde\mu:=\mu+\frac n{r_e}\,\mathcal L^n$ is a non-negative
distribution, therefore a non-negative Radon measure on $\Omega$, and
$\mu=\widetilde\mu-\frac n{r_e}\,\mathcal L^n$ is a locally finite signed Radon
measure. It remains to show that $\widetilde\mu(\Omega)<\infty$.

\emph{Step 3: volume of the boundary layer.}
Since $\partial\Omega$ is compact and $C^1$, it is covered by finitely many
charts in which it is the graph $x_n=g(x')$ of a Lipschitz function, and
there is $s_0>0$ such that every $x\in\Omega$ with $d(x)<s_0$ lies in the
cylinder of a chart containing a nearest boundary point $(y',g(y'))$ of $x$.
For such $x$,
\[
\abs{x_n-g(x')}\le\abs{x_n-g(y')}+\abs{g(y')-g(x')}
\le\bigl(1+\operatorname{Lip}g\bigr)\,d(x),
\]
so $\{0<d<s\}$ lies, in each chart, in a vertical strip of thickness
$2(1+\operatorname{Lip}g)s$ over a bounded base. Summing over the charts,
\begin{equation}\label{eq:layer}
\abs{\{0<d<s\}}\ \le\ c_\Omega\,s\qquad(0<s<s_0),
\end{equation}
with $c_\Omega$ depending only on the covering.

\emph{Step 4: total mass of $\widetilde\mu$.}
Fix $0<\sigma<s_0/2$ and, for $0<\eta<\sigma$, set
\[
\varphi_{\eta}:=\min\Bigl\{\frac{(d-\eta)^+}{\sigma},\,1\Bigr\},
\]
a Lipschitz function with values in $[0,1]$, supported in the compact set
$\{d\ge\eta\}\subset\Omega$, with $\nabla\varphi_\eta=\sigma^{-1}\nabla d$
on $\{\eta<d<\eta+\sigma\}$ and $\nabla\varphi_\eta=0$ almost everywhere
elsewhere. Since $d\in W^{1,\infty}(\Omega)$, the definition of $\mu=-\Delta d$
gives
\begin{equation}\label{eq:test-tilde-mu}
\int_\Omega\varphi\,d\widetilde\mu
=\int_\Omega\nabla d\cdot\nabla\varphi\,dx+\frac n{r_e}\int_\Omega\varphi\,dx
\qquad\text{for }\varphi\in C_c^\infty(\Omega),
\end{equation}
and \eqref{eq:test-tilde-mu} extends to $\varphi=\varphi_\eta$: the
mollifications $\varphi_{\eta,j}\in C_c^\infty(\Omega)$ of $\varphi_\eta$
converge to it uniformly and in $W^{1,1}(\Omega)$, with supports in a fixed
compact $K\Subset\Omega$, so the left-hand side passes to the limit because
$\widetilde\mu(K)<\infty$ and the right-hand side because
$\nabla d\in L^\infty(\Omega)$. Using $\abs{\nabla d}=1$ almost everywhere,
$0\le\varphi_\eta\le1$, and \eqref{eq:layer} with $\eta+\sigma<2\sigma<s_0$,
\[
\int_\Omega\varphi_{\eta}\,d\widetilde\mu
=\frac{\abs{\{\eta<d<\eta+\sigma\}}}{\sigma}+\frac n{r_e}\int_\Omega\varphi_\eta\,dx
\ \le\ \frac{c_\Omega(\eta+\sigma)}{\sigma}+\frac n{r_e}\abs{\Omega}
\ \le\ 2c_\Omega+\frac n{r_e}\abs{\Omega}.
\]
As $\eta\downarrow0$ the functions $\varphi_\eta$ increase pointwise to
$\min\{d/\sigma,1\}$, which equals $1$ on $\{d\ge\sigma\}$; by monotone
convergence,
\[
\widetilde\mu\bigl(\{d\ge\sigma\}\bigr)
\ \le\ \int_\Omega\min\{d/\sigma,1\}\,d\widetilde\mu
\ \le\ 2c_\Omega+\frac n{r_e}\abs{\Omega},
\]
independently of $\sigma$. Letting $\sigma\downarrow0$ gives
$\widetilde\mu(\Omega)\le2c_\Omega+n\abs{\Omega}/r_e<\infty$, and therefore
\[
\abs{\mu}(\Omega)\ \le\ \widetilde\mu(\Omega)+\frac n{r_e}\abs{\Omega}
\ \le\ 2c_\Omega+\frac{2n}{r_e}\abs{\Omega}\ <\ \infty . \qedhere
\]
\end{proof}

\end{document}
